\section{Universal driven closure for two factors}
\label{sec:two-factor}

Let $X\in\R^{n\times d_x}$ contain the data points as rows and consider
\begin{equation}
  f=XW^\top a,
  \qquad
  W\in\R^{m\times d_x},
  \qquad
  a\in\R^m.
  \label{eq:two-factor-model}
\end{equation}
Let $G=XX^\top$ be the fixed data Gram matrix and $r=\nabla_f\ell(f)$.

\begin{proposition}[Universal two-factor tangent-kernel closure]
\label[proposition]{prop:two-factor-closure}
For every hidden width $m$, input dimension $d_x$, finite dataset, differentiable loss, and parameter state,
\begin{equation}
  K=XW^\top WX^\top+\norm{a}^2G,
  \label{eq:two-factor-kernel}
\end{equation}
and gradient flow satisfies
\begin{equation}
  \boxed{
  \dot f=-Kr,
  \qquad
  \dot K=-(Gr)f^\top-f(Gr)^\top-2(f^\top r)G.}
  \label{eq:two-factor-kdot}
\end{equation}
Hence $(f,K)$ is an exact closed state.  The $K$-component depends only on $(f,r)$ and the fixed data geometry $G$, not on the hidden factors or directly on $K$.
\end{proposition}

\begin{proof}
For sample $i$,
\begin{equation}
  \nabla_W f_i=a x_i^\top,
  \qquad
  \nabla_a f_i=Wx_i,
\end{equation}
which gives \eqref{eq:two-factor-kernel}.  Put $s=X^\top r$.  Gradient flow is
\begin{equation}
  \dot W=-a s^\top,
  \qquad
  \dot a=-Ws.
  \label{eq:two-factor-parameter-flow}
\end{equation}
Because $G=XX^\top$ is fixed, $\dot G=0$.  Differentiating \eqref{eq:two-factor-kernel} gives
\begin{equation}
  \dot K
  =X(\dot W^\top W+W^\top\dot W)X^\top
   +\left(\frac{\dd}{\dd t}\norm{a}^2\right)G.
  \label{eq:two-factor-proof-step-one}
\end{equation}
Here the scalar derivative multiplying the fixed matrix $G$ is
\[
  \frac{\dd}{\dd t}\norm{a}^2
  =2\langle a,\dot a\rangle
  =-2a^\top Ws
  =-2f^\top r.
\]
Substituting \eqref{eq:two-factor-parameter-flow} into \eqref{eq:two-factor-proof-step-one} therefore yields
\begin{equation}
  \dot K
  =-(Xs)(a^\top WX^\top)-(XW^\top a)(Xs)^\top
   -2(a^\top Ws)G.
  \label{eq:two-factor-proof-substitution}
\end{equation}
Thus the last contribution is $\bigl(2\langle a,\dot a\rangle\bigr)G=-2(f^\top r)G$; it comes from differentiating $\norm{a}^2$, never from differentiating the fixed matrix $G$.  Since $Xs=Gr$ and $XW^\top a=f$, \eqref{eq:two-factor-kdot} follows.
\end{proof}

\subsection{Driven closure rather than a constitutive self-law}

Define
\begin{equation}
  B_G(f,r)=-(Gr)f^\top-f(Gr)^\top-2(f^\top r)G.
  \label{eq:BG}
\end{equation}
The exact reduced system is
\begin{equation}
  \dot f=-Kr(f),
  \qquad
  \dot K=B_G\bigl(f,r(f)\bigr).
  \label{eq:joint-driven-system}
\end{equation}
It is autonomous on the product state $(f,K)$, but it is not a constitutive field $K\mapsto\dot K$ on the positive-semidefinite cone.  The kernel has no instantaneous self-dynamics: its influence on its own future is indirect, through the prediction path that it helps generate.  In this precise sense the closure is \emph{driven}.  Consequently, the exact case in which the hidden representation disappears is also a case in which a kernel-only gradient-flow question becomes vacuous unless one freezes the prediction, treats it as a control, or works on the product space.

\subsection{A nontrivial foliation that the closure does not need}

\begin{proposition}[Balancing invariant and leaf independence]
\label[proposition]{prop:two-factor-invariant}
The matrix
\begin{equation}
  C=WW^\top-aa^\top
  \label{eq:two-factor-invariant}
\end{equation}
is conserved under \eqref{eq:two-factor-parameter-flow}.  Nevertheless, the closed law \eqref{eq:two-factor-kdot} is independent of $C$.
\end{proposition}

\begin{proof}
Using \eqref{eq:two-factor-parameter-flow},
\begin{align}
  \frac{\dd}{\dd t}(WW^\top)
  &=-a s^\top W^\top-Ws a^\top,\\
  \frac{\dd}{\dd t}(aa^\top)
  &=-Ws a^\top-a s^\top W^\top.
\end{align}
The derivatives are equal, so $\dot C=0$.  The formula \eqref{eq:two-factor-kdot} contains no $C$, proving leaf independence.
\end{proof}

This separates two notions that are often conflated.  An invariant foliation may exist without its labels being required by a reduced state.  The relevant question is not whether invariants exist, but whether the projected velocity changes across the leaves identified by the chosen observable.

\subsection{Exact elimination and a non-fading memory law}

The $K$-equation in \eqref{eq:joint-driven-system} is integrable once the prediction path is known.  Define
\begin{equation}
  \mathcal Q_G(f,u)
  =(Gu)f^\top+f(Gu)^\top+2(f^\top u)G
  =-B_G(f,u).
  \label{eq:QG}
\end{equation}

\begin{proposition}[Exact prediction-only Volterra equation]
\label[proposition]{prop:volterra}
Fix $K_0$ and write $r_t=r(f_t)$.  The Markovian system \eqref{eq:joint-driven-system} is equivalent to
\begin{equation}
  K_t=K_0-\int_0^t\mathcal Q_G(f_s,r_s)\,\dd s
  \label{eq:K-integral}
\end{equation}
and the closed prediction-only equation
\begin{equation}
  \boxed{
  \dot f_t=-K_0r_t
  +\int_0^t\mathcal Q_G(f_s,r_s)r_t\,\dd s.}
  \label{eq:volterra}
\end{equation}
Equivalently,
\begin{align}
  \dot f_t={}&-K_0r_t
  +\int_0^t\Bigl[
    (Gr_s)(f_s^\top r_t)
    +f_s(r_s^\top G r_t)\nonumber\\
  &\hspace{31mm}+2(f_s^\top r_s)Gr_t
  \Bigr] \,\dd s.
  \label{eq:volterra-expanded}
\end{align}
\end{proposition}

\begin{proof}
Integrating $\dot K=-\mathcal Q_G(f,r)$ gives \eqref{eq:K-integral}.  Substitution into $\dot f_t=-K_t r_t$ gives \eqref{eq:volterra}.  Conversely, reconstructing $K$ by \eqref{eq:K-integral} recovers both equations in \eqref{eq:joint-driven-system}.
\end{proof}

Equation \eqref{eq:volterra} contains the complete past $s\in[0,t]$ without an explicit lag-decay factor.  Old contributions may become small because the states $(f_s,r_s)$ become small, but the representation has no structural fading kernel that would justify truncating the history to a short window.  Conditional on $K_0$, there is no additional stochastic or orthogonal-dynamics remainder: the integral carries the entire effect of the eliminated kernel.  The same deterministic training dynamics are therefore Markovian in $(f,K)$ and history-dependent in $f$ alone.  Memory is a property of the chosen observable, not of the underlying vector field in isolation.  This is an elementary, completely resolved instance of the projection principle emphasized by Mori--Zwanzig theory \citep{zwanzig1961memory,mori1965transport,chorin2002memory}.
